\small
\renewcommand{\arraystretch}{1.3}

\begin{longtable}{|>{\raggedright\arraybackslash}p{0.13\linewidth}|>{\raggedright\arraybackslash}p{0.14\linewidth}|>{\raggedright\arraybackslash}p{0.14\linewidth}|>{\raggedright\arraybackslash}p{0.14\linewidth}|>{\raggedright\arraybackslash}p{0.18\linewidth}|>{\raggedright\arraybackslash}p{0.08\linewidth}|>{\raggedright\arraybackslash}p{0.07\linewidth}|}
	
	\caption{Failure Mode and Effect Analysis (FMEA)}
	\label{tab:fmea}\\
	\hline
	
	\textbf{Component} &
	\textbf{Failure Modes} &
	\textbf{Effects of Failure} &
	\textbf{Causes of Failure} &
	\textbf{Recommended Action} &
	\textbf{Safety Requirements} &
	\textbf{Ref.} \\
	\hline
	\endfirsthead
	
	\multicolumn{7}{c}{Table \thetable\ -- Continued}\\
	\hline
	
	\textbf{Component} &
	\textbf{Failure Modes} &
	\textbf{Effects of Failure} &
	\textbf{Causes of Failure} &
	\textbf{Recommended Action} &
	\textbf{Safety Requirements} &
	\textbf{Ref.} \\
	\hline
	\endhead
	
	\hline
	\multicolumn{7}{r}{Continued on next page}\\
	\endfoot
	
	\hline
	\endlastfoot
	
	Property Configuration &
	The system is unable to recognize the user's connected input device (such 
	as a controller), when selected in the feature variation menu &
	Users waste their time and can become confused/frustrated when attempting 
	to troubleshoot input device connection &
	A. System does not provide sufficient guidance to the user to connect their 
	input device
	
	B. System does not correctly monitor device interfaces to receive input 
	(Bluetooth, USB) &
	A. Provide detailed guidance to the user on how to connect each input 
	device; prompt the user to select the device they wish to connect 
	(PlayStation DualSense controller, Xbox Series X controller, etc.) and 
	provide guidance specific to that device
	
	B. Monitor both USB and Bluetooth interfaces for connection to each type of 
	supported input device when the user opts to connect a new device &
	CREATE NEW REQS FOR THIS &
	H1-A \\
	\hline
	
	World Navigation &
	User is unable to understand where their player character can move next in 
	the overall game world &
	User can become confused/frustrated, and waste gameplay time staying in the 
	main game world instead of progressing forwards &
	System does not provide adequate feedback to the user regarding where the 
	player character can move next; valid locations are indistinguishable from 
	invalid locations, or a possible path to take to reach the end of the game 
	is unclear &
	Highlight valid locations that the player character can move to in the game 
	world, and indicate invalid locations. Provide further guidance to the user 
	regarding a possible path forward if an extended period of time is spent 
	attempting to reach invalid locations &
	CREATE NEW REQS FOR THIS &
	H2-A \\
	\hline
	
	Combat System &
	The system displays flashing or intense graphics &
	Users with photosensitive epilepsy can potentially suffer an epileptic 
	seizure &
	The system displays excessive flashing graphics without warning users with 
	photosensitive epilepsy &
	Include a detailed flashing lights or epilepsy warning when the system 
	launches, and ensure the game does not display excessive flashing graphics &
	CREATE NEW REQS FOR THIS &
	H3-A \\
	\hline
	
	Progression and Rewards &
	Rewards provided to the user for succeeding in an enemy encounter are lost 
	when exiting the encounter to the game world &
	Gameplay progression is interrupted by earned rewards not being given to 
	the user, preventing them from using rewards to improve their in-game 
	capabilities; this can cause the user to become frustrated &
	Variables storing enemy encounter rewards are erroneously overwritten by 
	another value during transitions between game scenes &
	Retain and update enemy encounter reward variables (during encounters) in 
	an additional memory location separate from the combat system, to ensure a 
	backup value is available &
	CREATE NEW REQS FOR THIS &
	H4-A \\
	\hline
	
	Central Hub &
	Failure mode 1 &
	effect 1 &
	cause 1 &
	action1 &
	CREATE NEW REQS FOR THIS &
	H5-A \\
	\hline
	
	Progress Management &
	Failure mode 1 &
	effect 1 &
	cause 1 &
	action1 &
	CREATE NEW REQS FOR THIS &
	H6-A \\
	\hline
	
	User Interface &
	Failure mode 1 &
	effect 1 &
	cause 1 &
	action1 &
	CREATE NEW REQS FOR THIS &
	H7-A \\
	\hline
	
\end{longtable}